\begin{table}[tbp]
\centering
\caption{Predeclared primary evidence-tier rules and abstention outcomes.}
\label{tab:evidence-tier-rules}
\scriptsize
\begin{tabular}{p{0.39\linewidth}rp{0.40\linewidth}}
\toprule
Required diagnostic or condition & Threshold & Audit outcome \\
\midrule
Complete common-method comparison & required & Missing input $\Rightarrow$ inconclusive \\
Selection-aware nested interval & required & Unavailable $\Rightarrow$ inconclusive \\
Unique stable post-transition time placebos & 50 & Fewer dates $\Rightarrow$ inconclusive \\
Raw time-placebo probability & 0.10 & Above cutoff $\Rightarrow$ not supported \\
BH-adjusted time-placebo q value & 0.10 & Unavailable $\Rightarrow$ inconclusive; above cutoff $\Rightarrow$ not supported \\
Fixed-weight conditional diagnostic interval & exclude 0 & Otherwise $\Rightarrow$ not supported \\
Leave-one-donor-out direction fraction & 0.90 & Below cutoff $\Rightarrow$ not supported \\
All available required conditions & pass & Supported candidate discontinuity \\
\bottomrule
\end{tabular}
\par\smallskip
\begin{minipage}{0.94\linewidth}
\footnotesize\textit{Note.} The rules synthesize observational diagnostics; they are not a supervised classifier, instrument-fault label, or physical-causality test. The nested interval is required for tier assignment but is not claimed to have synthetic coverage calibration.
\end{minipage}
\end{table}
